\subsection{Deriving the formula for \(\varepsilon^{ij}\) and \(\mu^{ij}\)}
Using the tools from the previous section, we can express Maxwell's equations in a
coordinate-free manner.
Maxwell's equations
\begin{align*}
\nabla \cdot \boldsymbol{E} &= \frac{\rho}{\varepsilon_0}\\
\nabla \cdot \boldsymbol{H} &= 0\\
\nabla\times \boldsymbol{E} &=-\mu_0\frac{\partial \boldsymbol{H}}{\partial t}\\
\nabla\times \boldsymbol{H} &=\varepsilon_0\frac{\partial \boldsymbol{E}}{\partial t} + \boldsymbol{J}
\end{align*}
become
\begin{align}
\begin{split}\label{maxwellGeneral}
\frac{1}{\sqrt{g}}\partial_i\left(\sqrt{g}E^i\right) &= \frac{\rho}{\varepsilon_0}\\
\frac{1}{\sqrt{g}}\partial_i\left(\sqrt{g}H^i\right) &=0 \\
\epsilon^{ijk} \partial_jE_k &= -\mu_0\frac{\partial H^i}{\partial t}\\
\epsilon^{ijk} \partial_jH_k &= \varepsilon_0\frac{\partial E^i}{\partial t} + J^i.
\end{split}
\end{align}
Lowering the indices on \(H^i\) and \(E^i\), and writing the Levi-Civita tensor explicitly as
\(\epsilon^{ijk}=\pm\frac{1}{\sqrt{g}}[ijk]\), we obtain
\begin{align*}
\partial_i\left(\sqrt{g}g^{ij}E_j\right) &= \frac{\sqrt{g}\rho}{\varepsilon_0}\\
\partial_i\left(\sqrt{g}g^{ij}H_j\right) &=0\\
[ijk]\partial_jE_k &= -\mu_0\frac{\partial\left(\pm \sqrt{g}g^{ij}H_j\right)}{\partial t}\\
[ijk]\partial_jH_k &= \varepsilon_0\frac{\partial \left(\pm \sqrt{g}g^{ij}E_j\right)}{\partial t} \pm \sqrt{g}J^i.
\end{align*}
If we adopt the constitutive relations 
\begin{align}
D^i&=\varepsilon_0\varepsilon^{ij}E_j \nonumber\\
B^i&=\mu_0\mu^{ij}H_j \nonumber\\
\varepsilon^{ij}&=\mu^{ij}=\pm\sqrt{g}g^{ij}\label{cartesiantensors}
\end{align}
and define \(\varrho=\sqrt{g}\rho\) and \(\mathcal{J}^i=\pm\sqrt{g}J^i\), we
recover Maxwell's equations in macroscopic form:
\begin{align*}
\partial_i D^i&=\varrho\\
\partial_i B^i&=0\\
[ijk]\partial_jE_k&=-\frac{\partial B^i}{\partial t}\\
[ijk]\partial_jH_k&=\frac{\partial D^i}{\partial t}+\mathcal{J}^i.
\end{align*}
It is worth pausing to summarize our approach. We start with
Maxwell's equations in Cartesian coordinates, rewrite them in arbitrary coordinates,
and then reinterpret them in the original Cartesian system. Equation
\eqref{cartesiantensors} then provides the material parameters that induce the
desired transformation. The limitation of this approach is that the
transformation must be described in Cartesian coordinates. We cannot directly apply
formula \eqref{cartesiantensors} to transformation \eqref{transformation}
because it is expressed in cylindrical coordinates. To gain more flexibility, we start
with an arbitrary coordinate system, transform to a different arbitrary coordinate
system, and reinterpret the equations in the original coordinates. This is:
we start with general coordinate system with metric \(\gamma\):
\begin{align}\label{maxwellgamma}
\begin{split}
\frac{1}{\sqrt{\gamma}}\partial_i\left(\sqrt{\gamma}E^i\right) &= \frac{\rho}{\varepsilon_0}\\
\frac{1}{\sqrt{\gamma}}\partial_i\left(\sqrt{\gamma}H^i\right) &=0\\
[ijk] \partial_jE_k &= -\mu_0\frac{\partial \left(\sqrt{\gamma} H^i\right)}{\partial t}\\
[ijk] \partial_jH_k &= \varepsilon_0\frac{\partial \left(\sqrt{\gamma} E^i\right)}{\partial t} + \sqrt{\gamma}J^i,
\end{split}
\end{align}
then we go to a different general coordinate system, this is
\eqref{maxwellGeneral}, then we interpret \eqref{maxwellGeneral} in
the coordinates of \eqref{maxwellgamma},
which yields the constitutive relations:
\begin{align}
	D^i&=\varepsilon_0\varepsilon^{ij}E_j\nonumber\\
	B^i&=\mu_0\mu^{ij}H_j\nonumber\\
	\varepsilon^{ij}&=\mu^{ij}=\pm\frac{\sqrt{g}}{\sqrt{\gamma}}g^{ij}\label{tensorformulawithg}.
\end{align}
Equation \eqref{tensorformulawithg} provides a recipe for computing the material
parameters starting from an arbitrary coordinate system with an arbitrary
coordinate transformation. However, the inverse metric tensor that appears in
the formula makes it impractical, since this tensor is not immediately available
from the transformation. This tensor describes the geometry that light perceives, but we would
prefer a formula that uses the coordinate transformation and our
original metric tensor---that is, the metric of the space we started with. For this,
we apply the relationship between metric tensors from the previous
section, namely equation \eqref{tensorrelations}:
\begin{equation*}
	\boldsymbol{g'}=\Lambda^T \boldsymbol{g} \Lambda.
\end{equation*}
Taking the determinant of both sides yields
\[g'=g\left(\det\Lambda\right)^2.\]
Taking the inverse of both sides gives
\[
\boldsymbol{g}^{-1}=\Lambda(\boldsymbol{g}')^{-1}\Lambda^T.
\]
Thus
\begin{align*}
	\mu^{ij}=\varepsilon^{ij}&=\pm\frac{1}{\sqrt{\gamma}}\sqrt{g}g^{ij}\\
		&=\pm\frac{1}{\sqrt{\gamma}}\frac{\sqrt{g'}}{|\det\Lambda|}\Lambda(\boldsymbol{g}')^{-1}\Lambda^T\\
		&=\frac{\sqrt{g'}}{\sqrt{\gamma}}\frac{\Lambda(\boldsymbol{g}')^{-1}\Lambda^T}{\det\Lambda}.
\end{align*}
This formula is more practical because it depends on the coordinate
transformation \(\Lambda\) and \(\boldsymbol{g}'\), which represents the metric tensor
of the coordinate system we started with---in our particular case, cylindrical coordinates.

\subsection{Using the formula}
